\documentclass[a4paper,11pt]{article}
\usepackage[T2A]{fontenc}
\usepackage[utf8]{inputenc}
\usepackage[russian]{babel}
\usepackage{amsmath,amssymb,mathtools}
\usepackage{geometry}
\usepackage{microtype}
\usepackage{booktabs,tabularx,array}
\usepackage{enumitem}
\usepackage{tikz}
\usepackage{xcolor}
\usepackage{hyperref}
\usepackage{fancyhdr}
\usepackage{icomma}
\geometry{left=22mm,right=22mm,top=21mm,bottom=23mm}
\definecolor{accent}{HTML}{087D80}
\definecolor{muted}{HTML}{46515B}
\definecolor{error}{HTML}{A33A3A}
\definecolor{success}{HTML}{166B67}
\hypersetup{colorlinks=true,linkcolor=accent,urlcolor=accent,pdftitle={Ревью домашней работы — Николь Саркисьянц},pdfauthor={Лёвин Артём Александрович}}
\setlength{\parindent}{0pt}
\setlength{\parskip}{6pt}
\renewcommand{\arraystretch}{1.25}
\setlist[enumerate]{leftmargin=1.5em,itemsep=3pt,topsep=3pt}
\pagestyle{fancy}
\fancyhf{}
\fancyhead[L]{\footnotesize\color{muted}Математика · ревью домашней работы}
\fancyhead[R]{\footnotesize\color{muted}8 ноября 2026 г.}
\fancyfoot[C]{\scriptsize\color{muted}Лёвин Артём Александрович эксклюзивно для Николь Саркисьянц}
\fancyfoot[R]{\footnotesize\thepage}
\renewcommand{\headrulewidth}{0.25pt}
\renewcommand{\footrulewidth}{0pt}
\newcommand{\keyidea}{\textcolor{accent}{\textbf{Ключевая идея.}}\ }
\newcommand{\errorlabel}{\textcolor{error}{\textbf{Ошибка.}}\ }
\newcommand{\result}{\textcolor{success}{\textbf{Итог.}}\ }
\newcommand{\tasksection}[1]{\par\vspace{8pt}{\Large\color{accent}\bfseries Задание № #1}\par\vspace{2pt}\rule{\linewidth}{0.35pt}\par}
\begin{document}
\thispagestyle{empty}
\begin{center}
\vspace*{33mm}
{\large\color{accent}МАТЕМАТИКА}\par
\vspace{23mm}
{\huge Ревью домашней работы}\par
\vspace{11mm}
{\Large\color{accent}Кредит и тригонометрическое уравнение}\par
\vspace{28mm}
{\large Николь Саркисьянц}\par
\vspace{11mm}
{8 ноября 2026 г.}\par
\vfill
{\small Лёвин Артём Александрович}\par
{\small эксклюзивно для Николь Саркисьянц}\par
\vspace{15mm}
\begin{tikzpicture}[x=1cm,y=1cm]
\draw[accent,line width=0.8pt] (0,0) .. controls (1.1,0.45) and (2,-0.5) .. (3.2,0)
 .. controls (4.25,0.47) and (5.1,-0.48) .. (6.2,0)
 .. controls (7.1,0.4) and (8,-0.43) .. (9,0);
\end{tikzpicture}
\vspace{4mm}
\end{center}
\newpage
\section*{Сводка по заданиям с ошибками}
В таблицу включены только задания, в которых потребовались исправления. Нажмите на номер, чтобы перейти к подробному разбору.
\par\vspace{6pt}\noindent
{\small
\begin{tabularx}{\linewidth}{@{}>{\raggedright\arraybackslash}p{12mm} >{\raggedright\arraybackslash}p{26mm} X >{\raggedright\arraybackslash}p{31mm} >{\raggedright\arraybackslash}p{34mm}@{}}
\toprule
№ & Тема & Суть ошибки & Ваш ответ & Правильный ответ\\
\midrule
\hyperref[task:13]{\textbf{13}} & Кредит & Во втором году выбран неверный остаток долга и платёж. & 1\,324\,800 руб. & \textbf{1\,339\,200 руб.}\\
\addlinespace[6pt]
\hyperref[task:14]{\textbf{14}} & Тригонометрия & Неверны ограничение для корня и знак при преобразовании. & $0,\frac\pi3,\pi,$\newline$\frac{5\pi}3,2\pi$ & $0,\frac{2\pi}3,$\newline$\pi,2\pi$\\
\bottomrule
\end{tabularx}
}\par
\vspace{15pt}
\textcolor{accent}{\textbf{Что важно отработать}}\par
В задаче о кредите сопоставляйте остаток долга с условием каждого июля. В тригонометрическом уравнении отдельно контролируйте область определения квадратных корней и знак при использовании основного тригонометрического тождества.
\newpage
\thispagestyle{empty}
\vspace*{0.37\textheight}
\begin{center}
{\huge\color{accent}Разбор ошибок}\par
\vspace{11pt}
{\large Два задания · полные решения и упражнения}
\end{center}
\newpage
\phantomsection\label{task:13}
\tasksection{13}
\textbf{Условие.} В июле взяли кредит 900\,000 рублей на три года. Каждый январь остаток долга увеличивается на 20\%. К июлю первого и второго годов выплат остаток долга должен уменьшаться на 20\% относительно предыдущего июля. К июлю третьего года кредит погашен. Найдите общую сумму выплат.

\textbf{Ваше решение.} Использована таблица изменения долга. В первом году верно получены проценты 180\,000 рублей, долг после начисления 1\,080\,000 рублей и платёж 360\,000 рублей. На следующий июль правильно оставлено 720\,000 рублей. Во втором году после начисления 20\% верно получены 864\,000 рублей, однако вновь записан платёж 360\,000 рублей; отсюда появился остаток 504\,000 рублей. В третьем году найдено 604\,800 рублей, а общая сумма получилась
\[
360\,000\cdot2+604\,800=1\,324\,800.
\]
\textbf{Ваш ответ:} 1\,324\,800 рублей.

\errorlabel Во втором году не выполнено обязательное условие о снижении июльского остатка на 20\%. Оно относится к \emph{остатку долга}, а не к размеру платежа.

\textbf{Где именно возникает ошибка.} Последний правильный шаг:
\[
720\,000\cdot 1{,}2=864\,000.
\]
Следующий шаг в Вашей таблице: платёж второго года равен 360\,000 рублей, а остаток --- 504\,000 рублей. Арифметически вычитание верно, но полученный остаток противоречит условию: он должен составлять 80\% от 720\,000 рублей.

\textbf{Почему это неверно.} После первого года долг равен 720\,000 рублей. Значит, к июлю второго года он должен снизиться на $0{,}20\cdot720\,000=144\,000$ рублей по сравнению с прошлым июлем. Требуемый остаток равен $720\,000-144\,000=576\,000$ рублей. Платёж определяется разницей между долгом после январского начисления и необходимым остатком в июле.

\keyidea Удобно сначала найти все июльские остатки, а затем рассчитать платежи за каждый год.

\textbf{Исправление с последнего верного шага.}
\[
D_2=720\,000\cdot0{,}8=576\,000,
\qquad P_2=864\,000-576\,000=288\,000.
\]
Теперь в третьем году начисление 20\% происходит на 576\,000 рублей:
\[
P_3=576\,000\cdot1{,}2=691\,200.
\]
\pagebreak[4]
\textbf{Полное правильное решение.}
\begin{enumerate}
\item После первого июльского платежа должно остаться $900\,000\cdot0{,}8=720\,000$ рублей. Поэтому первый платёж:
\[
P_1=900\,000\cdot1{,}2-720\,000=360\,000.
\]
\item После второго июльского платежа должно остаться $720\,000\cdot0{,}8=576\,000$ рублей. Следовательно,
\[
P_2=720\,000\cdot1{,}2-576\,000=288\,000.
\]
\item В третьем году кредит погашают полностью. После январского начисления долг равен $576\,000\cdot1{,}2=691\,200$ рублей. Это и есть третий платёж.
\item Складываем все выплаты:
\[
S=360\,000+288\,000+691\,200
 =\boxed{1\,339\,200\text{ рублей}}.
\]
\end{enumerate}
\textbf{Проверка.} Уменьшение остатков: $900\,000\to720\,000\to576\,000\to0$. Начисленные проценты за три года равны $180\,000+144\,000+115\,200=439\,200$ рублей. Поэтому другая проверка общей суммы:
\[
900\,000+439\,200=1\,339\,200.
\]
\result Общая сумма выплат составляет $\boxed{1\,339\,200\text{ рублей}}$.

\textbf{Задача на закрепление.} В июле взяли кредит 600\,000 рублей на три года. Каждый январь остаток долга увеличивается на 10\%. К июлю первого и второго годов выплат остаток долга должен уменьшаться на 20\% относительно предыдущего июля. К июлю третьего года долг полностью погашен. Найдите общую сумму выплат.

\newpage
\phantomsection\label{task:14}
\tasksection{14}
\textbf{Условие.} Решите уравнение и найдите все его корни на отрезке $[0;2\pi]$:
\[
\sqrt{2\cos^3x-\sin^2x-2\cos x+\sin x}
=\sqrt{\cos\left(\frac\pi2-x\right)}.
\]
\textbf{Ваше решение.} Правая часть приведена к $\sqrt{\sin x}$, далее при сравнении подкоренных выражений получено
\[
2\cos x(\cos^2x-1)-\sin^2x=0.
\]
Затем в следующей строке записано
\[
2\cos x\sin^2x-\sin^2x=0,
\]
что приводит к уравнению $\sin^2x(2\cos x-1)=0$. На листе также выписано ограничение $\sin x\le0$. После отбора на отрезке записаны значения $0$, $\pi$, $2\pi$, $\frac\pi3$, $\frac{5\pi}3$.

\textbf{Ваш ответ:} $0$, $\frac\pi3$, $\pi$, $\frac{5\pi}3$, $2\pi$.

\errorlabel Здесь две независимые ошибки.
\begin{enumerate}
\item \textbf{Область определения.} Для $\sqrt{\sin x}$ необходимо $\sin x\ge0$. В записи указано обратное неравенство. Подкоренное выражение слева также должно быть \emph{неотрицательным}; нуль допустим.
\item \textbf{Знак при замене.} Из тождества $\cos^2x-1=-\sin^2x$ следует
\[
2\cos x(\cos^2x-1)=-2\cos x\sin^2x.
\]
В Вашем переходе знак «минус» перед первым слагаемым потерян. Из-за этого появилось неверное уравнение $2\cos x-1=0$.
\end{enumerate}
\textbf{Последний правильный шаг и первый неправильный.} Правильно:
\[
2\cos x(\cos^2x-1)-\sin^2x=0.
\]
Неверно:
\[
2\cos x\sin^2x-\sin^2x=0.
\]
\keyidea Квадратные корни равны тогда и только тогда, когда их подкоренные выражения равны и оба неотрицательны. После упрощения обязательно возвращайтесь к ограничениям.

\textbf{Исправление с последнего правильного шага.}
\begin{align*}
2\cos x(\cos^2x-1)-\sin^2x&=0,\\
-2\cos x\sin^2x-\sin^2x&=0,\\
-\sin^2x(2\cos x+1)&=0.
\end{align*}
Отсюда $\sin x=0$ или $\cos x=-\frac12$.

\textbf{Полное правильное решение.}
\begin{enumerate}
\item Формула приведения даёт
\[
\cos\left(\frac\pi2-x\right)=\sin x.
\]
Поэтому исходное уравнение равносильно системе
\[
\begin{cases}
2\cos^3x-\sin^2x-2\cos x+\sin x=\sin x,\\
\sin x\ge0.
\end{cases}
\]
\item Вычитаем $\sin x$ и выносим общий множитель:
\begin{align*}
2\cos^3x-\sin^2x-2\cos x&=0,\\
2\cos x(\cos^2x-1)-\sin^2x&=0,\\
-2\cos x\sin^2x-\sin^2x&=0,\\
\sin^2x(2\cos x+1)&=0.
\end{align*}
\item Решаем два простых уравнения:
\[
\sin x=0\Rightarrow x=\pi n,\ n\in\mathbb Z;
\qquad \cos x=-\frac12.
\]
На $[0;2\pi]$ первая ветвь даёт $0$, $\pi$, $2\pi$. Вторая даёт $\frac{2\pi}3$ и $\frac{4\pi}3$.
\item Проверяем $\sin x\ge0$: для $0$, $\pi$, $2\pi$ синус равен нулю, для $\frac{2\pi}3$ положителен, а для $\frac{4\pi}3$ отрицателен. Последнее значение исключаем.
\end{enumerate}
\textbf{Проверка.} Для четырёх оставшихся значений выполнено $\sin x\ge0$ и равенство подкоренных выражений. При $\frac{2\pi}3$ имеем $\cos x=-\frac12$, поэтому
\[
2\cos^3x-\sin^2x-2\cos x+\sin x=\sin x.
\]
При $0$, $\pi$, $2\pi$ обе части исходного уравнения равны нулю. Значения $\frac\pi3$ и $\frac{5\pi}3$ из Вашего ответа не подходят: при первом левое подкоренное выражение отрицательно, а при втором отрицателен $\sin x$.

\result На отрезке $[0;2\pi]$ получаем
\[
\boxed{x\in\left\{0,\ \frac{2\pi}{3},\ \pi,\ 2\pi\right\}}.
\]
\textbf{Задача на закрепление.} Решите уравнение и найдите все его корни на отрезке $[0;2\pi]$:
\[
\sqrt{2\sin^3x-\cos^2x-2\sin x+\cos x}
=\sqrt{\sin\left(\frac\pi2-x\right)}.
\]
\end{document}
